\chapter[29]{Compact Two-Dimensional Ancient Solutions}
\label{ch:chow-iv-compact-two-dimensional-ancient-solutions}
\dcref{ch29}
\source{chow-et-al-ricci-flow-part-iv:page-0090}

Let \(g(t)=v(t)^{-1}g_{S^2}\), \(t\in(-\infty,0)\), be a maximal ancient
Ricci flow on the two-sphere.  The positive function \(v\) is called the
pressure.  In cylindrical and stereographic coordinates we write the same
metric as
\[
 g(t)=\widehat v(t)^{-1}g_{\mathrm{cyl}}
     =\overline v(t)^{-1}g_{\mathrm{euc}}.
\]

\section{Classification and the model solution}

\begin{theorem}[Classification of ancient solutions on the two-sphere]
\label{thm:chow-iv-classification-ancient-two-sphere}
\source{chow-et-al-ricci-flow-part-iv:page-0090}
\uses{prop:chow-iv-q-vanishes-backward, prop:chow-iv-q-zero-characterizes-king-rosenau}
\dcref{ch29:29.1}
\group{Compact Ancient Solutions}


Every compact simply connected ancient solution to the Ricci flow in
dimension two is either a round shrinking two-sphere or the rotationally
symmetric King--Rosenau solution.
\end{theorem}

\begin{theorem}[Two-dimensional Type I and Type II ancient solutions]
\label{thm:chow-iv-two-dimensional-type-i-ii}
\source{chow-et-al-ricci-flow-part-iv:page-0094}
\dcref{ch29:29.2}
\group{Compact Ancient Solutions}

Let \(g(t)\) be a compact ancient Ricci flow on \(S^2\).
\begin{enumerate}
\item If \(\sup_{S^2\times(-\infty,-1]}|t|R<\infty\), then \(g(t)\) is a
round shrinking sphere.
\item If this supremum is infinite, backward blow-ups at suitable points and
times converge to a cigar soliton.
\end{enumerate}
\end{theorem}

\begin{lemma}[Characterization of the King--Rosenau solution]
\label{lem:chow-iv-characterization-king-rosenau}
\source{chow-et-al-ricci-flow-part-iv:page-0095}
\dcref{ch29:29.3}
\group{King--Rosenau Solution}

Suppose an ancient solution has pressure
\(v(\psi,t)=\alpha(t)+\beta(t)\cos^2\psi\), where \(\alpha,\beta>0\) and
\(\alpha(t)\to\infty\) as \(t\nearrow0\).  Then, for some \(\mu>0\),
\[
 \alpha(t)=-\mu\operatorname{csch}(4\mu t),\qquad
 \beta(t)=-\mu\tanh(2\mu t),
\]
and the solution is the King--Rosenau solution.
\end{lemma}

\begin{lemma}[Cheeger--Gromov convergence to a cigar at the poles]
\label{lem:chow-iv-king-rosenau-pole-cigar-limit}
\source{chow-et-al-ricci-flow-part-iv:page-0096,chow-et-al-ricci-flow-part-iv:page-0097}
\dcref{ch29:29.4}
\group{King--Rosenau Solution}

For any \(t_i\to-\infty\), conformal diffeomorphisms fixing the two poles and
normalizing the metric at either pole pull the time-translated King--Rosenau
solutions back to a sequence converging in pointed smooth Cheeger--Gromov
sense to a cigar soliton.
\end{lemma}

\begin{lemma}[Pointed backward limits of the King--Rosenau solution]
\label{lem:chow-iv-king-rosenau-pointed-limits}
\source{chow-et-al-ricci-flow-part-iv:page-0097}
\dcref{ch29:29.5}
\group{King--Rosenau Solution}

For base points \(((s_i,\theta_i),t_i)\) in cylindrical coordinates with
\(t_i\to-\infty\), the pointed limit is a cigar, a flat cylinder, or a cigar
with the opposite orientation according as \(2t_i-s_i\) tends to
\(-\infty\), a finite number, or \(+\infty\).
\end{lemma}

\section{Pressure estimates and backward limits}

\begin{lemma}[Uniform Sobolev and Holder bounds for the pressure]
\label{lem:chow-iv-pressure-sobolev-holder}
\source{chow-et-al-ricci-flow-part-iv:page-0098}
\dcref{ch29:29.6}
\group{Pressure Estimates}

For every \(p<\infty\) and \(\alpha\in(0,1)\), uniformly for \(t\le-1\),
\[
 \|v(t)\|_{W^{2,p}(S^2)}\le C(p),\qquad
 \|v(t)\|_{C^{1,\alpha}(S^2)}\le C(\alpha).
\]
\end{lemma}

\begin{remark}[Pressure and the soliton potential]
\label{rem:chow-iv-pressure-soliton-potential}
\source{chow-et-al-ricci-flow-part-iv:page-0098}
\dcref{ch29:29.7}
\group{Pressure Estimates}

The identities \(\Delta_g\log v=R-2v\) and
\(g=e^{-\log v}g_{S^2}\) make \(\log v\) analogous to a steady-soliton
potential.
\end{remark}

\begin{lemma}[Regularity of the squared pressure gradient]
\label{lem:chow-iv-pressure-gradient-regularity}
\source{chow-et-al-ricci-flow-part-iv:page-0099}
\uses{lem:chow-iv-pressure-sobolev-holder}
\dcref{ch29:29.8}
\group{Pressure Estimates}


For every \(p<\infty\) and \(\alpha\in(0,1)\), uniformly for \(t\le-1\),
\[
 \bigl\||\nabla v|^2\bigr\|_{W^{2,p}}\le C(p),\qquad
 \bigl\||\nabla v|^2\bigr\|_{C^{1,\alpha}}\le C(\alpha).
\]
\end{lemma}

\begin{remark}[Intrinsic curvature derivative bounds]
\label{rem:chow-iv-intrinsic-curvature-derivatives}
\source{chow-et-al-ricci-flow-part-iv:page-0099}
\dcref{ch29:29.9}
\group{Pressure Estimates}

The curvature derivative estimates give
\(|\nabla_g^mR|_g\le C_m\) for \(t\le-1\), but do not directly bound the
derivatives measured with the round metric.
\end{remark}

\begin{lemma}[Weighted higher pressure estimates]
\label{lem:chow-iv-weighted-higher-pressure-estimates}
\source{chow-et-al-ricci-flow-part-iv:page-0099}
\uses{lem:chow-iv-pressure-gradient-regularity}
\dcref{ch29:29.10}
\group{Pressure Estimates}


For every \(\alpha\in(0,1)\), uniformly for \(t\le-1\),
\[
 \|v^{3/2}\|_{C^{2,\alpha}}+
 \|\sqrt v\,\nabla^2v\|_{C^\alpha}+
 \|v^2\|_{C^{3,\alpha}}\le C(\alpha).
\]
\end{lemma}

\begin{lemma}[Regularity of the weighted Hessian]
\label{lem:chow-iv-weighted-pressure-hessian}
\source{chow-et-al-ricci-flow-part-iv:page-0099,chow-et-al-ricci-flow-part-iv:page-0100}
\uses{lem:chow-iv-weighted-higher-pressure-estimates}
\dcref{ch29:29.11}
\group{Pressure Estimates}


For \(p<\infty\) and \(\alpha\in(0,1)\), uniformly for \(t\le-1\),
\[
 \|v\nabla^2v\|_{W^{1,p}}\le C(p),\qquad
 \|v\nabla^2v\|_{C^\alpha}+\|v\Delta_{S^2}v\|_{C^\alpha}\le C(\alpha).
\]
\end{lemma}

\begin{lemma}[Convergence to the backward pressure]
\label{lem:chow-iv-backward-pressure-convergence}
\source{chow-et-al-ricci-flow-part-iv:page-0102}
\uses{lem:chow-iv-pressure-gradient-regularity, lem:chow-iv-weighted-higher-pressure-estimates, lem:chow-iv-weighted-pressure-hessian}
\dcref{ch29:29.12}
\group{Backward Pressure Limit}


The limit \(v_\infty=\lim_{t\to-\infty}v(t)\) lies in
\(C^{1,\alpha}\) for all \(\alpha\in(0,1)\).  The convergence is in
\(C^{1,\alpha}\) globally, in \(C^{3,\alpha}\) locally on
\(\{v_\infty>0\}\), and weakly in \(W^{2,2}\); moreover
\(|\nabla v(t)|^2\to|\nabla v_\infty|^2\) in \(C^{1,\alpha}\).
\end{lemma}

\begin{definition}[Weak stationary pressure]
\label{def:chow-iv-weak-stationary-pressure}
\source{chow-et-al-ricci-flow-part-iv:page-0103}
\dcref{ch29:29.13}
\group{Backward Pressure Limit}

A function \(\widetilde v\in W^{1,2}(S^2)\) is a weak solution of the
stationary pressure equation if
\[
 \int_{S^2}\bigl(-\widetilde v\langle\nabla\widetilde v,\nabla\varphi\rangle
 -|\nabla\widetilde v|^2\varphi+2\widetilde v^2\varphi\bigr)
 \,d\mu_{S^2}=0
\]
for every smooth test function \(\varphi\).
\end{definition}

\begin{lemma}[The backward pressure is stationary]
\label{lem:chow-iv-backward-pressure-stationary}
\source{chow-et-al-ricci-flow-part-iv:page-0103}
\uses{lem:chow-iv-backward-pressure-convergence, def:chow-iv-weak-stationary-pressure}
\dcref{ch29:29.14}
\group{Backward Pressure Limit}


The limit \(v_\infty\) is a weak solution of
\[
 v_\infty\Delta_{S^2}v_\infty-|\nabla v_\infty|^2+2v_\infty^2=0.
\]
It is smooth on \(\{v_\infty>0\}\), where
\(v_\infty^{-1}g_{S^2}\) is flat.
\end{lemma}

\section{Isoperimetry and the round case}

\begin{proposition}[Infinite expansion implies a round two-sphere]
\label{prop:chow-iv-zero-backward-pressure-round}
\source{chow-et-al-ricci-flow-part-iv:page-0104,chow-et-al-ricci-flow-part-iv:page-0118,chow-et-al-ricci-flow-part-iv:page-0119}
\uses{lem:chow-iv-isoperimetric-rigidity, lem:chow-iv-radius-linear-time}
\dcref{ch29:29.16}
\group{Round Case}


If \(v_\infty\equiv0\), then \(g(t)\) is a round shrinking two-sphere.
\end{proposition}

\begin{example}[Isoperimetric constant of the unit two-sphere]
\label{ex:chow-iv-unit-sphere-isoperimetric-constant}
\source{chow-et-al-ricci-flow-part-iv:page-0104}
\dcref{ch29:29.17}
\group{Isoperimetry}

For the unit round sphere, every latitude circle has isoperimetric ratio
\(4\pi\), and every embedded loop has ratio at least \(4\pi\).  Thus the
isoperimetric constant is \(4\pi\).
\end{example}

\begin{lemma}[Isoperimetric constant of metrics on \(S^2\)]
\label{lem:chow-iv-isoperimetric-minimizers}
\source{chow-et-al-ricci-flow-part-iv:page-0104,chow-et-al-ricci-flow-part-iv:page-0105,chow-et-al-ricci-flow-part-iv:page-0106,chow-et-al-ricci-flow-part-iv:page-0107}
\dcref{ch29:29.18}
\group{Isoperimetry}

For every smooth metric \(\widetilde g\) on \(S^2\), its isoperimetric
constant satisfies \(I(\widetilde g)\le4\pi\) and depends continuously on
\(\widetilde g\) in the \(C^0\) topology.  If \(I(\widetilde g)<4\pi\), the
infimum is attained by a connected smooth embedded loop of constant geodesic
curvature.
\end{lemma}

\begin{lemma}[Rigidity of the optimal isoperimetric constant]
\label{lem:chow-iv-isoperimetric-rigidity}
\source{chow-et-al-ricci-flow-part-iv:page-0107,chow-et-al-ricci-flow-part-iv:page-0108}
\uses{lem:chow-iv-isoperimetric-minimizers}
\dcref{ch29:29.19}
\group{Isoperimetry}


If \(I(\widetilde g)=4\pi\), then \(\widetilde g\) is round.  If a Ricci
flow on \(S^2\) has isoperimetric constant \(4\pi\) at one time, then it is a
round shrinking solution.
\end{lemma}

\begin{lemma}[Evolution of isoperimetric ratios of parallel curves]
\label{lem:chow-iv-isoperimetric-ratio-evolution}
\source{chow-et-al-ricci-flow-part-iv:page-0108}
\dcref{ch29:29.20}
\group{Isoperimetry}

For equidistant curves \(\gamma_\rho\) about an embedded loop at time \(t_0\),
their isoperimetric ratios satisfy
\[
 \left(\partial_t-\partial_\rho^2-
 \frac{\int_{\gamma_\rho}\kappa_\rho\,ds}{L(\gamma_\rho)}\partial_\rho\right)
 \log I(\rho,t)\big|_{t=t_0}
 \ge \frac{4\pi-I(\rho,t_0)}{4\pi|t_0|}.
\]
\end{lemma}

\begin{proposition}[Decay of the isoperimetric constant]
\label{prop:chow-iv-isoperimetric-decay}
\source{chow-et-al-ricci-flow-part-iv:page-0109}
\uses{lem:chow-iv-isoperimetric-ratio-evolution, lem:chow-iv-isoperimetric-rigidity}
\dcref{ch29:29.21}
\group{Isoperimetry}


If the ancient solution is not round, then \(I(t)<4\pi\),
\(I(t)\to0\) as \(t\to-\infty\), and \(I(t)\) obeys the quantitative
backward decay estimate stated in (29.64)--(29.65).
\end{proposition}

\begin{proposition}[Uniform bounds for minimizing loops]
\label{prop:chow-iv-isoperimetric-loop-bounds}
\source{chow-et-al-ricci-flow-part-iv:page-0109,chow-et-al-ricci-flow-part-iv:page-0110}
\uses{prop:chow-iv-isoperimetric-decay, lem:chow-iv-isoperimetric-minimizers}
\dcref{ch29:29.22}
\group{Isoperimetry}


For every minimizing loop \(\gamma_t\) and \(t\le-1\), there is \(C\ge1\)
such that
\[
 C^{-1}\le L_{g(t)}(\gamma_t)\le C.
\]
\end{proposition}

\begin{lemma}[Comparable areas of the isoperimetric regions]
\label{lem:chow-iv-isoperimetric-region-areas}
\source{chow-et-al-ricci-flow-part-iv:page-0111}
\uses{prop:chow-iv-isoperimetric-loop-bounds}
\dcref{ch29:29.23}
\group{Isoperimetry}


There is \(c>0\) such that the two regions cut out by any minimizing loop
satisfy
\[
 c|t|\le A_\pm(\gamma_t)\le8\pi|t|\qquad(t\le-1).
\]
\end{lemma}

\begin{lemma}[Existence of an asymptotic cylinder]
\label{lem:chow-iv-asymptotic-cylinder-minimizing-loop}
\source{chow-et-al-ricci-flow-part-iv:page-0111,chow-et-al-ricci-flow-part-iv:page-0112}
\uses{prop:chow-iv-isoperimetric-loop-bounds, lem:chow-iv-isoperimetric-region-areas}
\dcref{ch29:29.24}
\group{Backward Geometry}


For \(t_i\to-\infty\), minimizers \(\gamma_{t_i}\), and
\(p_i\in\gamma_{t_i}\), a subsequence of
\((S^2,g(t+t_i),p_i)\) converges smoothly to a static flat cylinder.  The
loops converge to an embedded geodesic cross-section, so sufficiently old
minimizing loops are centers of arbitrarily good necks.
\end{lemma}

\begin{lemma}[Connectedness of distance circles]
\label{lem:chow-iv-connected-distance-circles}
\source{chow-et-al-ricci-flow-part-iv:page-0113,chow-et-al-ricci-flow-part-iv:page-0114}
\dcref{ch29:29.25}
\group{Backward Geometry}

There is \(\varepsilon>0\) such that, in a positively curved two-sphere, if
an \(\varepsilon\)-neck separates a point \(p\) from one end, then the first
distance circle from \(p\) meeting the neck is connected.
\end{lemma}

\begin{lemma}[Distance circles near the center of a neck]
\label{lem:chow-iv-distance-circles-near-neck-center}
\source{chow-et-al-ricci-flow-part-iv:page-0114,chow-et-al-ricci-flow-part-iv:page-0115}
\dcref{ch29:29.26}
\group{Backward Geometry}

For every \(\delta>0\), sufficiently good neck coordinates make the relevant
distance circles \(C^0\)-close to cylindrical cross-sections and make their
lengths \(\delta\)-close to the cylindrical circumference.
\end{lemma}

\begin{lemma}[Area lower bound for balls]
\label{lem:chow-iv-area-lower-bound-balls}
\source{chow-et-al-ricci-flow-part-iv:page-0115,chow-et-al-ricci-flow-part-iv:page-0116}
\dcref{ch29:29.27}
\group{Backward Geometry}

There is \(c_0>0\) such that for every \(p\in S^2\), every \(r\ge3\), and
every \(t\le-1\),
\[
 \operatorname{Area}_{g(t)}B_{g(t)}(p,r)\ge c_0r.
\]
\end{lemma}

\begin{lemma}[Area bounds in terms of time]
\label{lem:chow-iv-polar-ball-area-time}
\source{chow-et-al-ricci-flow-part-iv:page-0116}
\uses{lem:chow-iv-connected-distance-circles, lem:chow-iv-isoperimetric-region-areas}
\dcref{ch29:29.28}
\group{Backward Geometry}


Let \(x_t^\pm\) maximize distance from \(\gamma_t\) in the two complementary
regions and write that distance as \(\rho_t^\pm\).  For all sufficiently
negative \(t\),
\[
 c_0|t|\le\operatorname{Area}_{g(t)}B(x_t^\pm,\rho_t^\pm)
 \le8\pi|t|.
\]
\end{lemma}

\begin{lemma}[Linear-time bounds for polar radii]
\label{lem:chow-iv-radius-linear-time}
\source{chow-et-al-ricci-flow-part-iv:page-0117}
\uses{lem:chow-iv-area-lower-bound-balls, lem:chow-iv-polar-ball-area-time}
\dcref{ch29:29.29}
\group{Backward Geometry}


There is \(c_1>0\) such that
\[
 c_1^{-1}|t|\le\rho_t^\pm\le c_1|t|\qquad(t\le-1).
\]
In particular, \(\operatorname{diam}(S^2,g(t))\le C|t|\).
\end{lemma}

\begin{proposition}[The plane is not a backward pointwise limit]
\label{prop:chow-iv-plane-not-backward-pointwise-limit}
\source{chow-et-al-ricci-flow-part-iv:page-0119,chow-et-al-ricci-flow-part-iv:page-0120,chow-et-al-ricci-flow-part-iv:page-0121}
\uses{lem:chow-iv-distance-circles-near-neck-center, lem:chow-iv-radius-linear-time}
\dcref{ch29:29.30}
\group{Backward Geometry}


No fixed point \(q\in S^2\) can have
\((S^2,g(t_i+t),q)\) converge without rescaling to the flat plane as
\(t_i\to-\infty\).
\end{proposition}

\begin{lemma}[Nonround solutions cannot have plane limits]
\label{lem:chow-iv-nonround-no-plane-limit}
\source{chow-et-al-ricci-flow-part-iv:page-0121}
\uses{prop:chow-iv-plane-not-backward-pointwise-limit}
\dcref{ch29:29.31}
\group{Backward Geometry}


If \(g(t)\) is nonround and \(q_i\in S^2\), then no unrescaled pointed
smooth limit of \((S^2,g(t_i+t),q_i)\), \(t_i\to-\infty\), is the flat
plane.
\end{lemma}

\section{Classification of the backward pressure}

\begin{lemma}[Integral bound for positive circular averages]
\label{lem:chow-iv-positive-circular-average-bound}
\source{chow-et-al-ricci-flow-part-iv:page-0122}
\dcref{ch29:29.32}
\group{Backward Pressure Limit}

If the backward Euclidean conformal factor is finite at the origin, then for
every \(r_0>0\) and \(t\le-1\), the integral of the positive part of the
circular average of \(\log\overline u\) over \(B(0,r_0)\) is bounded by a
constant depending only on \(r_0\).
\end{lemma}

\begin{lemma}[Unbounded conformal factor concentrates curvature]
\label{lem:chow-iv-unbounded-conformal-factor-curvature-concentration}
\source{chow-et-al-ricci-flow-part-iv:page-0123,chow-et-al-ricci-flow-part-iv:page-0124}
\uses{lem:chow-iv-positive-circular-average-bound}
\dcref{ch29:29.33}
\group{Backward Pressure Limit}


If \(\overline u(z_0,t_i)\to\infty\), then every fixed Euclidean ball about
\(z_0\) contains asymptotically at least \(4\pi\) of total scalar curvature.
\end{lemma}

\begin{remark}[Round-sphere concentration]
\label{rem:chow-iv-round-sphere-curvature-concentration}
\source{chow-et-al-ricci-flow-part-iv:page-0123}
\dcref{ch29:29.34}
\group{Backward Pressure Limit}

For a round shrinking sphere the Euclidean conformal factor diverges at every
fixed point as \(t\to-\infty\), and every fixed disk captures asymptotically
\(4\pi\) of scalar curvature.
\end{remark}

\begin{lemma}[A nontrivial backward pressure has at most two zeros]
\label{lem:chow-iv-backward-pressure-at-most-two-zeros}
\source{chow-et-al-ricci-flow-part-iv:page-0124}
\uses{lem:chow-iv-unbounded-conformal-factor-curvature-concentration}
\dcref{ch29:29.35}
\group{Backward Pressure Limit}


Either \(v_\infty\equiv0\), or \(v_\infty\) has at most two zeros on \(S^2\).
\end{lemma}

\begin{proposition}[Infinite expansion or cylinder limit]
\label{prop:chow-iv-backward-pressure-cylinder}
\source{chow-et-al-ricci-flow-part-iv:page-0125,chow-et-al-ricci-flow-part-iv:page-0126}
\uses{prop:chow-iv-zero-backward-pressure-round, lem:chow-iv-backward-pressure-at-most-two-zeros, prop:chow-iv-plane-not-backward-pointwise-limit}
\dcref{ch29:29.36}
\group{Backward Pressure Limit}


For a nonround solution, after a conformal diffeomorphism, there is \(\mu>0\)
such that
\[
 v(t)\longrightarrow v_\infty(\psi,\theta)=\mu\cos^2\psi
\]
in \(C^{1,\alpha}(S^2)\) and smoothly on compact subsets away from the poles.
Equivalently, the backward metric limit on the punctured sphere is a flat
cylinder.
\end{proposition}

\begin{corollary}[Angular pressure-gradient bound]
\label{cor:chow-iv-angular-pressure-gradient-bound}
\source{chow-et-al-ricci-flow-part-iv:page-0126}
\uses{prop:chow-iv-backward-pressure-cylinder}
\dcref{ch29:29.37}
\group{Backward Pressure Limit}


For every nonround solution there is \(C<\infty\) such that
\[
 \left|\frac{\partial}{\partial\theta}\log v\right|\le C
 \qquad\text{on }S^2\times(-\infty,-1].
\]
\end{corollary}

\begin{proposition}[Backward convergence at the poles to cigar solitons]
\label{prop:chow-iv-poles-converge-to-cigars}
\source{chow-et-al-ricci-flow-part-iv:page-0127,chow-et-al-ricci-flow-part-iv:page-0128,chow-et-al-ricci-flow-part-iv:page-0129}
\uses{prop:chow-iv-backward-pressure-cylinder, lem:chow-iv-nonround-no-plane-limit, lem:chow-iv-rescaled-pressure-growth}
\dcref{ch29:29.38}
\group{Backward Geometry}


For every nonround ancient solution and every sequence \(t_i\to-\infty\), a
subsequence of the solutions normalized at either pole converges smoothly on
compact subsets to a cigar soliton.  Equivalently, suitable conformal
pullbacks converge in pointed smooth Cheeger--Gromov sense to a cigar.
\end{proposition}

\begin{lemma}[Growth of the rescaled pressure]
\label{lem:chow-iv-rescaled-pressure-growth}
\source{chow-et-al-ricci-flow-part-iv:page-0128}
\dcref{ch29:29.39}
\group{Backward Geometry}

For the pole-normalized Euclidean pressures \(\overline v_i\), there is a
uniform constant \(C\) such that
\[
 -C\le\log\overline v_i(r,\theta,0)\le C(1+r^2)
 \qquad\text{on }\mathbb R^2.
\]
\end{lemma}

\section{\texorpdfstring{The invariant \(Q\)}{The invariant Q}}

Let \(B=\operatorname{Sym}(\nabla^3v)\), let
\(\operatorname{TF}(B)\) be its totally trace-free part, and set
\[
 Q=v|\operatorname{TF}(B)|^2.
\]

\begin{lemma}[Norm of the trace-free symmetrized third derivative]
\label{lem:chow-iv-trace-free-third-derivative}
\source{chow-et-al-ricci-flow-part-iv:page-0130}
\dcref{ch29:29.40}
\group{The \(Q\) Invariant}

The tensor \(B=\operatorname{Sym}(\nabla^3v)\) satisfies
\[
 |\operatorname{TF}(B)|^2=|\nabla^3v|^2-\frac34|\nabla\Delta v|^2
 -|\nabla v|^2-\langle\nabla\Delta v,\nabla v\rangle.
\]
\end{lemma}

\begin{remark}[Irreducible decomposition of
\(\nabla^3v\)]
\label{rem:chow-iv-third-derivative-decomposition}
\source{chow-et-al-ricci-flow-part-iv:page-0130,chow-et-al-ricci-flow-part-iv:page-0131}
\dcref{ch29:29.42}
\group{The \(Q\) Invariant}

The irreducible decomposition of \(\nabla^3v\) gives the same trace-free
tensor as the symmetrized construction and makes the identity in Lemma 29.40
orthogonal.
\end{remark}

\begin{proposition}[\(Q\) is a heat subsolution]
\label{prop:chow-iv-q-heat-subsolution}
\source{chow-et-al-ricci-flow-part-iv:page-0132,chow-et-al-ricci-flow-part-iv:page-0133,chow-et-al-ricci-flow-part-iv:page-0134,chow-et-al-ricci-flow-part-iv:page-0135,chow-et-al-ricci-flow-part-iv:page-0136,chow-et-al-ricci-flow-part-iv:page-0137,chow-et-al-ricci-flow-part-iv:page-0138}
\uses{lem:chow-iv-symmetrized-third-derivative-evolution}
\dcref{ch29:29.44}
\group{The \(Q\) Invariant}


The quantity \(Q=v|\operatorname{TF}(B)|^2\) satisfies
\[
 (\partial_t-v\Delta_{S^2})Q\le0.
\]
Equivalently, because \(\Delta_g=v\Delta_{S^2}\), it is a subsolution of
the heat equation of \(g(t)\).
\end{proposition}

\begin{lemma}[Evolution of the symmetrized third derivative]
\label{lem:chow-iv-symmetrized-third-derivative-evolution}
\source{chow-et-al-ricci-flow-part-iv:page-0133}
\dcref{ch29:29.45}
\group{The \(Q\) Invariant}

The symmetric tensor \(B=\operatorname{Sym}(\nabla^3v)\) satisfies the
heat-type equation (29.145), obtained by differentiating the pressure
equation and commuting covariant derivatives on the round sphere.
\end{lemma}

\begin{proposition}[Vanishing of \(Q\) characterizes King--Rosenau]
\label{prop:chow-iv-q-zero-characterizes-king-rosenau}
\source{chow-et-al-ricci-flow-part-iv:page-0138,chow-et-al-ricci-flow-part-iv:page-0139,chow-et-al-ricci-flow-part-iv:page-0140,chow-et-al-ricci-flow-part-iv:page-0141,chow-et-al-ricci-flow-part-iv:page-0142,chow-et-al-ricci-flow-part-iv:page-0143,chow-et-al-ricci-flow-part-iv:page-0144,chow-et-al-ricci-flow-part-iv:page-0145}
\uses{lem:chow-iv-equivalence-q-euclidean-q, lem:chow-iv-characterization-king-rosenau, prop:chow-iv-backward-pressure-cylinder}
\dcref{ch29:29.46}
\group{The \(Q\) Invariant}


The identity \(Q\equiv0\) holds if and only if the ancient solution is round
or is the King--Rosenau solution.
\end{proposition}

\begin{example}[Euclidean expressions for the \(Q\)-tensor]
\label{ex:chow-iv-euclidean-q-expressions}
\source{chow-et-al-ricci-flow-part-iv:page-0139}
\dcref{ch29:29.47}
\group{The \(Q\) Invariant}

For a Euclidean pressure \(\widetilde v\), let
\[
 \widetilde a_{ijk}=\widetilde v_{ijk}-\frac14
 \bigl((\Delta\widetilde v)_i\delta_{jk}
 +(\Delta\widetilde v)_j\delta_{ik}
 +(\Delta\widetilde v)_k\delta_{ij}\bigr).
\]
Then \(|\widetilde a|^2\) is one quarter of the sum of the squares of
\(\widetilde v_{111}-3\widetilde v_{221}\) and
\(\widetilde v_{222}-3\widetilde v_{112}\); it vanishes for the
King--Rosenau pressure.
\end{example}

\begin{lemma}[Evolution of the Euclidean invariant]
\label{lem:chow-iv-euclidean-q-evolution}
\source{chow-et-al-ricci-flow-part-iv:page-0145,chow-et-al-ricci-flow-part-iv:page-0146}
\dcref{ch29:29.48}
\group{The \(Q\) Invariant}

For \(\overline Q=\overline v|\overline a|^2\) on \(\mathbb R^2\), the
evolution is the nonpositive heat-type identity (29.196).  In particular,
\(\overline Q\) is a subsolution of the Euclidean pressure heat operator.
\end{lemma}

\begin{proposition}[The invariant vanishes backward]
\label{prop:chow-iv-q-vanishes-backward}
\source{chow-et-al-ricci-flow-part-iv:page-0146,chow-et-al-ricci-flow-part-iv:page-0147,chow-et-al-ricci-flow-part-iv:page-0148,chow-et-al-ricci-flow-part-iv:page-0149,chow-et-al-ricci-flow-part-iv:page-0150}
\uses{prop:chow-iv-q-heat-subsolution, prop:chow-iv-poles-converge-to-cigars, lem:chow-iv-equivalence-q-euclidean-q, prop:chow-iv-backward-pressure-cylinder}
\dcref{ch29:29.49}
\group{The \(Q\) Invariant}


If the constant \(\mu\) in the backward cylindrical limit is positive, then
\(Q_{\max}(t)=\max_{S^2}Q(\cdot,t)\) is nonincreasing and
\[
 \lim_{t\to-\infty}Q_{\max}(t)=0.
\]
Consequently \(Q\equiv0\).
\end{proposition}

\begin{lemma}[Equivalence of spherical and Euclidean \(Q\)]
\label{lem:chow-iv-equivalence-q-euclidean-q}
\source{chow-et-al-ricci-flow-part-iv:page-0150,chow-et-al-ricci-flow-part-iv:page-0151,chow-et-al-ricci-flow-part-iv:page-0152,chow-et-al-ricci-flow-part-iv:page-0153}
\dcref{ch29:29.50}
\group{The \(Q\) Invariant}

Let \(\sigma:S^2\setminus\{S\}\to\mathbb R^2\) be stereographic
projection, and let \(Q\) and \(\overline Q\) be the spherical and Euclidean
invariants associated to \(v\) and \(\overline v\).  Then
\[
 Q\circ\sigma^{-1}=\overline Q.
\]
\end{lemma}

\begin{remark}[Problem 29.15: completeness of the positive set]
\label{rem:chow-iv-completeness-positive-pressure-set}
\dcref{ch29:29.15}
\group{Open Problems}
\source{chow-et-al-ricci-flow-part-iv:page-0103}
Is \((\{v_\infty>0\},v_\infty^{-1}g_{S^2})\) directly seen to be complete?
The later classification shows that every nonempty such positive set is a
flat cylinder.
\end{remark}

\begin{remark}[Exercise 29.41: an alternative formula for \(Q\)]
\label{rem:chow-iv-q-alternative-formula}
\dcref{ch29:29.41}
\group{The \(Q\) Invariant}
\source{chow-et-al-ricci-flow-part-iv:page-0130}
If \(p=\nabla^2v+vg_{S^2}\) and \(P=\operatorname{tr}_g p\), then
\[
 \operatorname{TF}(B)=
 \operatorname{Sym}\!\left(\nabla\left(p-\frac34Pg_{S^2}\right)\right),
\]
and \(\nabla_i p_{jk}\) is totally symmetric.
\end{remark}

\begin{remark}[Exercise 29.43: vanishing on King--Rosenau]
\label{rem:chow-iv-q-vanishes-king-rosenau}
\dcref{ch29:29.43}
\group{The \(Q\) Invariant}
\source{chow-et-al-ricci-flow-part-iv:page-0132}
For the King--Rosenau pressure, the displayed Hessian and third-derivative
identities imply \(\operatorname{TF}(B)=0\), hence \(Q=0\).
\end{remark}
